% Cote 119: offprints, reports, courses, CV and letters (typed and printed part): archive cote 119, site page 224 (PDF page 225).
% typescript. Draft compared: none.
% AI-transcribed by Claude Opus 5.5 (claude-opus-5-5, adaptive thinking, effort high), Message Batches API, batch r119_full (msgbatch_01DwZNm5Y7YsKBWyAN4sFs8w), 2026-11-13; not yet checked by eye.
\scanpage{119/224}
\mspage{. 11 .}

les relations entre schémas abéliens et groupes de
Barsotti-Tate associés sont explorées et exploitées
également dans [17] et dans [16, SGA 7, Exp. IX].

\uline{Références} : J. Tate, B. Mazur, A. Néron, L. Illusie,
J.N. Katz, W. Messing, I. Manin.

9. \uline{Arithmétique} ([16 SGA 5, Exp XVI][18, Exp III])

Ma contribution principale a consisté (en
collaboration avec M. Artin) en la démonstration de
la rationalité des fonctions $L$ associées à des
faisceaux $\ell$-adiques généraux sur des variétés algébriques
sur des corps finis, comprenant comme cas particulier
les fonctions $L$ associées à des caractères de groupes
finis opérant sur de telles variétés. S'inspirant des
conjectures de Weil, on arrive en effet à exprimer ces
fonctions $L$ en termes de produits alternés de poly-%
nômes caractéristiques de Frobenius opérant sur la
``cohomologie à support propre'' de la variété envisagée.
Bien au delà d'une simple question de rationalité, ces
résultats ouvrent la voie à une \uline{approche} cohomologique
systématique d'invariants arithmétiques subtils comme
les fonctions $\zeta$ et $L$ des variétés, et l'interprétation
en termes arithmétiques de théorèmes tels que les
théorèmes de dualité (démontrés à l'heure actuelle) et
de Lefschetz pour les sections hyperplanes (non démontrés
encore en car. $> 0$) . Il y a là un champ d'étude immense,
qui par la nature des choses devrait se trouver, tôt
ou tard, centré sur la notion de ``motif'' (dénominateur
commun des divers types de cohomologie qu'on sait attacher
à une variété algébrique) - mais qui probablement ne
sera jamais exploré jusqu'au bout, l'heure de ce genre
d'investigations étant déjà passée (même si rares sont
ceux qui en ont pris conscience).

\margin{224}
